\documentclass{article}
\usepackage{amsmath}
\usepackage{amssymb}
\usepackage{siunitx}

\title{KdV Equation Theory}
\author{}
\date{}

\begin{document}

\section{The Korteweg-de Vries Equation}

The KdV equation describes weakly nonlinear long waves in dispersive media:

\begin{equation}
u_t + 6uu_x + u_{xxx} = 0
\end{equation}

where:
\begin{itemize}
    \item $u(x,t)$ -- wave amplitude
    \item $u_t = \frac{\partial u}{\partial t}$ -- temporal derivative
    \item $u_x = \frac{\partial u}{\partial x}$ -- spatial derivative
    \item $u_{xxx} = \frac{\partial^3 u}{\partial x^3}$ -- third spatial derivative
\end{itemize}

\section{Soliton Solutions}

The KdV equation admits exact soliton solutions -- localized waves that maintain their shape.

\subsection{Single Soliton}

\begin{equation}
u(x,t) = 2\kappa^2 \sech^2\left(\kappa(x - 4\kappa^2 t)\right)
\end{equation}

Parameters:
\begin{itemize}
    \item $\kappa$ -- soliton parameter (controls amplitude and width)
    \item Amplitude: $A = 2\kappa^2$
    \item Width: $w \sim 1/\kappa$
    \item Speed: $c = 4\kappa^2$
    \item Center: $x_c(t) = 4\kappa^2 t$
\end{itemize}

\subsection{Multi-Soliton Solutions}

For $N$ solitons with parameters $\kappa_i$:

\begin{equation}
u(x,t) = 2 \frac{\partial^2}{\partial x^2} \ln \det\left(\delta_{ij} + \frac{2\kappa_i\kappa_j}{\kappa_i^2 + \kappa_j^2} e^{-(\kappa_i + \kappa_j)x + 4(\kappa_i^3 + \kappa_j^3)t}\right)
\end{equation}

\section{Pseudo-Spectral Method}

\subsection{Discretization}

Periodic domain $[0, L)$ with $N$ equally spaced points:

\begin{align}
x_j &= j \Delta x, \quad j = 0, 1, \dots, N-1 \\
\Delta x &= \frac{L}{N}
\end{align}

\subsection{Fourier Representation}

\begin{align}
u_j(t) &= \frac{1}{N} \sum_{k=0}^{N-1} \hat{u}_k(t) e^{i k \theta_j} \\
\theta_j &= \frac{2\pi j}{N}
\end{align}

\subsection{Derivatives in Fourier Space}

\begin{align}
\widehat{u_x}_k &= i k \hat{u}_k \\
\widehat{u_{xxx}}_k &= (ik)^3 \hat{u}_k = -i k^3 \hat{u}_k
\end{align}

\section{Time Integration: ETD-RK4}

Exponential Time Differencing with Runge-Kutta 4:

\begin{align}
\hat{u}_t &= A\hat{u} + N(u) \\
\text{where } A\hat{u} &= -i k^3 \hat{u} \text{ (linear term)} \\
N(u) &= -\mathcal{F}\{6uu_x\} \text{ (nonlinear term in Fourier space)}
\end{align}

ETD-RK4 steps:

\begin{align}
\hat{u}_1 &= e^{\frac{h}{2}A}\hat{u}_n + \frac{h}{2}e^{\frac{h}{2}A}N(u_n) \\
\hat{u}_2 &= e^{\frac{h}{2}A}\hat{u}_n + \frac{h}{2}e^{\frac{h}{2}A}N(u_1) \\
\hat{u}_3 &= e^{hA}\hat{u}_n + he^{hA}N(u_2) \\
\hat{u}_{n+1} &= e^{hA}\hat{u}_n + \frac{h}{6}\left(e^{hA}N(u_0) + 2e^{\frac{h}{2}A}N(u_1) + 2e^{\frac{h}{2}A}N(u_2) + N(u_3)\right)
\end{align}

where $u_j = \mathcal{F}^{-1}\{\hat{u}_j\}$.

\section{Error Analysis}

\subsection{Truncation Error}
- Spatial: Spectral (exponential for smooth solutions)
- Temporal: $O(\Delta t^4)$ for RK4

\subsection{Gibbs Phenomenon}
Not applicable for periodic smooth solutions.

\subsection{Aliasing}
Use 2/3 rule for dealiasing nonlinear terms:
\begin{equation}
\hat{u}_k = 0 \text{ for } |k| > \frac{2N}{3}
\end{equation}

\section{Conservation Laws}

KdV has infinitely many conserved quantities:

\begin{align}
\text{Mass: } &M = \int_0^L u \, dx \\
\text{Momentum: } &P = \int_0^L u^2 \, dx \\
\text{Energy: } &E = \int_0^L \left(u_x^2 - 2u^3\right) \, dx
\end{align}

\end{document}
